%------------------------------------------------------------------------------%
\documentclass[fleqn,utf8,aspectratio=169,12pt]{beamer}

\usepackage{calculo_varias_variaveis-1}

\title{Cilindros e Superfícies Quádricas}

%------------------------------------------------------------------------------%
\begin{document}

\addtitle

%------------------------------------------------------------------------------%
\section{Cilindros}

%------------------------------------------------------------------------------%
\begin{frame}{Cilindro Sólido}
  O \emph{cilindro sólido} é um corpo geométrico tridimensional \\
  com duas bases circulares paralelas e congruentes \\
  conectadas por uma superfície lateral curva
\end{frame}

\framedgraphic*{Cilindro Sólido}{cilindro_solido}

%------------------------------------------------------------------------------%
\begin{frame}{Superfície Cilindrica}
  O \emph{cilindro} é uma superfície gerada
  por meio de o movimento de uma linha reta \\
  ao longo de uma curva plana dada
  enquanto mantida paralela a uma reta fixada

  A curva é chamada de \emph{curva geradora}
\end{frame}

\framedgraphic*{Cilindros}{cilindro_superficie_1}
\framedgraphic*{Cilindros}{cilindro_superficie_2}
\framedgraphic*{Cilindros}{cilindro_superficie_3}
\framedgraphic*{Cilindros}{cilindro_superficie_4}

%------------------------------------------------------------------------------%
\section{Quádricas}

%------------------------------------------------------------------------------%
\begin{frame}{Cônicas e Quádricas}
  \emph{Cônica}:
  pontos no plano, $\R^2$, que são solução de uma equação de segundo grau
  \[
    Ax^{\emph{2}} + By^{\emph{2}} + Cxy + Dx + Ey + F = 0
  \]

  \pause
  \emph{Quádrica}:
  pontos no espaço, $\R^3$, que são solução de uma equação de segundo grau
  \[
    Ax^{\emph{2}} + By^{\emph{2}} + Cz^{\emph{2}} +
    Dxy  + Exz  + Fyz +
    Gx   + Hy   + Iz + G = 0
  \]
\end{frame}

%------------------------------------------------------------------------------%
\begin{frame}{Quádricas}
  \begin{itemize}
    \item Elipsoides       \\[5mm]
    \item Paraboloides     \\[5mm]
    \item Hiperboloides    \\[5mm]
    \item Cones elípticos  \\[5mm]
    \item Cilindro elípticos
  \end{itemize}
\end{frame}

%------------------------------------------------------------------------------%
\begin{frame}{Elipsoide}
\begin{columns}
\column{0.4\linewidth}
  Elipsoide canônico
  \pause
  \bigskip

  Centrado na origem e\\
  \pause
  alinhado com os eixos
  \pause
  \medskip
  {\large
  \[
    \frac{x^2}{a^2} +
    \frac{y^2}{b^2} +
    \frac{z^2}{c^2} = 1
  \]}
  \medskip
  \(a>0\), \(b>0\) e \(c>0\)
\column{0.6\linewidth}
  \pause
  \includegraphics{quadricas_elipsoide}
\end{columns}
\end{frame}

%------------------------------------------------------------------------------%
\begin{frame}{Elipsoide}
\begin{columns}
\column{0.32\linewidth}
  \(
    x=0
    \pause
    \qquad
    \frac{y^2}{b^2} +
    \frac{z^2}{c^2} = 1
  \)
  \pause

  \bigskip
  \includegraphics{corte_elipsoide_x0}
  \pause
\column{0.32\linewidth}
  \(
    y=0
    \pause
    \qquad
    \frac{x^2}{a^2} +
    \frac{z^2}{c^2} = 1
  \)
  \pause

  \bigskip
  \includegraphics{corte_elipsoide_y0}
  \pause
\column{0.32\linewidth}
  \(
    z=0
    \pause
    \qquad
    \frac{x^2}{a^2} +
    \frac{y^2}{b^2} = 1
  \)
  \pause

  \bigskip
  \includegraphics{corte_elipsoide_z0}
\end{columns}
\end{frame}

%------------------------------------------------------------------------------%
\begin{frame}{Paraboloide Elíptico}
\begin{columns}
\column{0.4\linewidth}
  Paraboloide elíptico canônico
  \pause
  \bigskip

  ``Centrado'' na origem e\\
  \pause
  alinhado com os eixos
  \pause
  \medskip
  {\large
  \[
    z = \frac{x^2}{a^2}
      + \frac{y^2}{b^2}
  \]}
  \pause
  \medskip
  \(a>0\) e \(b>0\)
\column{0.6\linewidth}
  \pause
  \includegraphics{quadricas_paraboloide_eliptico}
\end{columns}
\end{frame}

%------------------------------------------------------------------------------%
\begin{frame}{Paraboloide Elíptico}
\begin{columns}
\column{0.32\linewidth}
  \(
    x=0
    \pause
    \qquad
    z = \frac{y^2}{b^2}
  \)
  \pause

  \bigskip
  \includegraphics{corte_paraboloide_eliptico_x0}
  \pause
\column{0.32\linewidth}
  \(
    y=0
    \pause
    \qquad
    z = \frac{x^2}{a^2}
  \)
  \pause

  \bigskip
  \includegraphics{corte_paraboloide_eliptico_y0}
  \pause
\column{0.32\linewidth}
  \(
    z=1
    \pause
    \qquad
    \frac{x^2}{a^2} +
    \frac{y^2}{b^2} = 1
  \)
  \pause

  \bigskip
  \includegraphics{corte_paraboloide_eliptico_z1}
\end{columns}
\end{frame}

%------------------------------------------------------------------------------%
\begin{frame}{Paraboloide Hiperbólico}
\begin{columns}
\column{0.4\linewidth}
  Paraboloide hiperbólico canônico
  \pause
  \bigskip

  ``Centrado'' na origem e\\
  \pause
  alinhado com os eixos
  \pause
  \medskip
  {\large
  \[
    \frac{y^2}{b^2} -
    \frac{x^2}{a^2} =
    \frac{z}{c}
  \]}
  \pause
  \medskip
  \(a>0\), \(b>0\) e \(c>0\)
\column{0.6\linewidth}
  \pause
  \includegraphics{quadricas_paraboloide_hiperbolico}
\end{columns}
\end{frame}

%------------------------------------------------------------------------------%
\begin{frame}{Paraboloide Hiperbólico}
\begin{columns}
\column{0.32\linewidth}
  \(
    x=0
    \pause
    \qquad
    \frac{y^2}{b^2} =
    \frac{z}{c}
  \)
  \pause

  \bigskip
  \includegraphics{corte_paraboloide_hiperbolico_x0}
  \pause
\column{0.32\linewidth}
  \(
    y=0
    \pause
    \qquad
    \frac{x^2}{a^2} =
    -\frac{z}{c}
  \)
  \pause

  \bigskip
  \includegraphics{corte_paraboloide_hiperbolico_y0}
  \pause
\column{0.32\linewidth}
  \(
    z=0
    \pause
    \qquad
    \frac{y^2}{b^2} =
    \frac{x^2}{a^2}
  \)
  \pause

  \bigskip
  \includegraphics{corte_paraboloide_hiperbolico_z0}
\end{columns}
\end{frame}

%------------------------------------------------------------------------------%
\begin{frame}{Hiperboloide de Uma Folha}
\begin{columns}
\column{0.4\linewidth}
  Hiperboloide de uma folha canônico
  \pause
  \bigskip

  ``Centrado'' na origem e\\
  \pause
  alinhado com os eixos
  \pause
  \medskip
  {\large
  \[
    \frac{x^2}{a^2} +
    \frac{y^2}{b^2} -
    \frac{z^2}{c^2} = 1
  \]}
  \pause
  \medskip
  \(a>0\), \(b>0\) e \(c>0\)
\column{0.6\linewidth}
  \pause
  \includegraphics{quadricas_hiperboloide_uma_folha}
\end{columns}
\end{frame}

%------------------------------------------------------------------------------%
\begin{frame}{Hiperboloide de Uma Folha}
\begin{columns}
\column{0.32\linewidth}
  \(
    x=0
    \pause
    \qquad
    \frac{y^2}{b^2} -
    \frac{z^2}{c^2} = 1
  \)
  \pause

  \bigskip
  \includegraphics{corte_hiperboloide_uma_folha_x0}
  \pause
\column{0.32\linewidth}
  \(
    y=0
    \pause
    \qquad
    \frac{x^2}{a^2} -
    \frac{z^2}{c^2} = 1
  \)
  \pause

  \bigskip
  \includegraphics{corte_hiperboloide_uma_folha_y0}
  \pause
\column{0.32\linewidth}
  \(
    z=0
    \pause
    \qquad
    \frac{x^2}{a^2} +
    \frac{y^2}{b^2} = 1
  \)
  \pause

  \bigskip
  \includegraphics{corte_hiperboloide_uma_folha_z0}
\end{columns}
\end{frame}

%------------------------------------------------------------------------------%
\begin{frame}{Hiperbolóide de Duas Folhas}
\begin{columns}
\column{0.4\linewidth}
  Hiperboloide de duas folhas canônico
  \pause
  \bigskip

  ``Centrado'' na origem e\\
  \pause
  alinhado com os eixos
  \pause
  \medskip
  {\large
  \[
    - \frac{x^2}{a^2}
    - \frac{y^2}{b^2}
    + \frac{z^2}{c^2}
    = 1
  \]}
  \pause
  \medskip
  \(a>0\), \(b>0\) e \(c>0\)
\column{0.6\linewidth}
  \pause
  \includegraphics{quadricas_hiperboloide_duas_folhas}
\end{columns}
\end{frame}

%------------------------------------------------------------------------------%
\begin{frame}{Hiperboloide de Duas Folhas}
\begin{columns}
\column{0.32\linewidth}
  \(
    x=0
    \pause
    \qquad
    \frac{z^2}{c^2} -
    \frac{y^2}{b^2} = 1
  \)
  \pause

  \bigskip
  \includegraphics{corte_hiperboloide_duas_folhas_x0}
  \pause
\column{0.32\linewidth}
  \(
    y=0
    \pause
    \qquad
    \frac{z^2}{c^2} -
    \frac{x^2}{a^2} = 1
  \)
  \pause

  \bigskip
  \includegraphics{corte_hiperboloide_duas_folhas_y0}
  \pause
\column{0.32\linewidth}
  \(
    z=0
    \pause
    \qquad
    - \frac{x^2}{a^2}
    - \frac{y^2}{b^2} = 1
  \)
  \pause

  \bigskip
  \includegraphics{corte_hiperboloide_duas_folhas_z0}
\end{columns}
\end{frame}

%------------------------------------------------------------------------------%
\begin{frame}{Cone Elíptico}
\begin{columns}
\column{0.4\linewidth}
  Cone elíptico canônico
  \pause
  \bigskip

  Centrado na origem e\\
  \pause
  alinhado com os eixos
  \pause
  \medskip
  {\large
  \[
    \frac{x^{2}}{a^2} +
    \frac{y^{2}}{b^2} =
    \frac{z^{2}}{c^2}
  \]}
  \pause
  \medskip
  \(a>0\), \(b>0\) e \(c>0\)
\column{0.6\linewidth}
  \pause
  \includegraphics{quadricas_cone_eliptico}
\end{columns}
\end{frame}

%------------------------------------------------------------------------------%
\begin{frame}{Cone Elíptico}
\begin{columns}
\column{0.32\linewidth}
  \(
    x=0
    \pause
    \qquad
    \frac{y^{2}}{b^2} =
    \frac{z^{2}}{c^2}
  \)
  \pause

  \bigskip
  \includegraphics{corte_cone_eliptico_x0}
  \pause
\column{0.32\linewidth}
  \(
    y=0
    \pause
    \qquad
    \frac{x^{2}}{a^2} =
    \frac{z^{2}}{c^2}
  \)
  \pause

  \bigskip
  \includegraphics{corte_cone_eliptico_y0}
  \pause
\column{0.32\linewidth}
  \(
    z=0
    \pause
    \qquad
    \frac{x^{2}}{a^2} +
    \frac{y^{2}}{b^2} = 0
  \)
  \pause

  \bigskip
  \includegraphics{corte_cone_eliptico_z0}
\end{columns}
\end{frame}

%------------------------------------------------------------------------------%
\begin{frame}{Cilindro Elíptico}
\begin{columns}
\column{0.4\linewidth}
  Cilindro elíptico canônico
  \pause
  \bigskip

  Centrado na origem e\\
  \pause
  alinhado com os eixos
  \pause
  \medskip
  {\large
  \[
    \frac{x^2}{a^2} +
    \frac{y^2}{b^2} = 1
  \]}
  \pause
  \medskip
  \(a>0\) e \(b>0\)
\column{0.6\linewidth}
  \pause
  \includegraphics{quadricas_cilindro_eliptico}
\end{columns}
\end{frame}

%------------------------------------------------------------------------------%
\begin{frame}{Cilindro Elíptico}
\begin{columns}
\column{0.32\linewidth}
  \(
    x=0
    \pause
    \qquad
    \frac{y^{2}}{b^2} = 1
  \)
  \pause

  \bigskip
  \includegraphics{corte_cilindro_eliptico_x0}
  \pause
\column{0.32\linewidth}
  \(
    y=0
    \pause
    \qquad
    \frac{x^{2}}{a^2} = 1
  \)
  \pause

  \bigskip
  \includegraphics{corte_cilindro_eliptico_y0}
  \pause
\column{0.32\linewidth}
  \(
    z=0
    \pause
    \qquad
    \frac{x^{2}}{a^2} +
    \frac{y^{2}}{b^2} = 1
  \)
  \pause

  \bigskip
  \includegraphics{corte_cilindro_eliptico_z0}
\end{columns}
\end{frame}

%------------------------------------------------------------------------------%
\section{Exemplos}

%------------------------------------------------------------------------------%
%------------------------------------------------------------------------------%
\begin{question}
  Encontre e esboce as curvas que representam o corte da superfície
  \[
    4x^2 + y^2 = 16
  \]
  pelos planos \(z=0\) e \(y=0\)
\end{question}

%------------------------------------------------------------------------------%
\begin{resolution}{Solução}
  No plano \(z=0\) a equação
  \[
    4x^2+y^2=16
  \]
  \pause
  permanece inalterada,
  \pause
  porém, agora representa a elipse
  \[
    \frac{x^2}{2^2} + \frac{y^2}{4^2} = 1
  \]
  \pause
  que intercepta o eixo $x$ nos pontos \((2, 0)\) e \((-2, 0)\)

  e o eixo $y$ nos pontos \((4, 0)\) e \((-4, 0)\)
\end{resolution}

%------------------------------------------------------------------------------%
\begin{resolution}{Solução}
  \onslide<+->{No plano \(y=0\) a equação
  \[
    4x^2+y^2=16
  \]}
  \onslide<+->{se reduz a}
  \begin{align*}
    \onslide<+->{4x^2 + 0 & = 16    \\}
    \onslide<+->{x^2      & = 4     \\}
    \onslide<+->{\abs{x}  & = 2     \\}
    \onslide<+->{x        & = \pm 2}
  \end{align*}
  \onslide<+->{que representa as retas verticais \(x=2\) e \(x=-2\)}
\end{resolution}

%------------------------------------------------------------------------------%
\begin{resolution}{Solução}
\begin{columns}
\column{0.5\linewidth}
  \centering
  \(z=0\) \\
  \medskip
  \includegraphics{C_04_exemplo_1_z0}
\column{0.5\linewidth}
  \centering
  \(y=0\) \\
  \medskip
  \includegraphics{C_04_exemplo_1_y0}
\end{columns}
\end{resolution}

%------------------------------------------------------------------------------%

%------------------------------------------------------------------------------%
\begin{question}
  Encontre e esboce as curvas que representam o corte da superfície
  \[
    x^2 + y^2 = 4 - z^2
  \]
  pelos planos \(z=0\) e \(y=0\)
\end{question}

%------------------------------------------------------------------------------%
\begin{resolution}{Solução}
  \onslide<+->{No plano \(z=0\) a equação
  \[
    x^2 + y^2 = 4 - z^2
  \]}
  \onslide<+->{se reduz a}
  \begin{align*}
    \onslide<+->{x^2 + y^2 & = 4 - 0 \\}
    \onslide<+->{x^2 + y^2 & = 2^2}
  \end{align*}
  \onslide<+->{que corresponde circunferência de raio 2 centrada na origem}
\end{question}

%------------------------------------------------------------------------------%
\begin{resolution}{Solução}
  \onslide<+->{No plano \(y=0\) a equação
  \[
    x^2 + y^2 = 4 - z^2
  \]}
  \onslide<+->{se reduz a}
  \begin{align*}
    \onslide<+->{x^2 + 0   & = 4 - z^2 \\}
    \onslide<+->{x^2 + z^2 & = 2^2}
  \end{align*}
  \onslide<+->{que corresponde circunferência de raio 2 centrada na origem}
\end{question}

%------------------------------------------------------------------------------%
\begin{resolution}{Solução}
\begin{columns}
  \column{0.5\linewidth}
  \centering
  \(z=0\) \\
  \medskip
  \includegraphics{C_04_exemplo_2_z0}
\column{0.5\linewidth}
  \centering
  \(y=0\) \\
  \medskip
  \includegraphics{C_04_exemplo_2_y0}
\end{columns}
\end{resolution}

%------------------------------------------------------------------------------%

%------------------------------------------------------------------------------%
\begin{question}
  Encontre e esboce as curvas que representam os cortes da superfície
  \[
    y^2 - x^2 = z
  \]
  pelos planos \(z=0\) e \(y=0\)
\end{question}

%------------------------------------------------------------------------------%
\begin{resolution}{Solução}
  \onslide<+->{No plano \(z=0\) a equação
  \[
    y^2 - x^2 = z
  \]}
  \onslide<+->{se reduz a}
  \begin{align*}
    \onslide<+->{y^2-x^2 & = 0       \\}
    \onslide<+->{y^2     & = x^2     \\}
    \onslide<+->{\abs{y} & = \abs{x} \\}
    \onslide<+->{y       & = \pm x}
  \end{align*}
  \onslide<+->{que corresponde as retas \(y=x\) e \(y=-x\)}
\end{resolution}

%------------------------------------------------------------------------------%
\begin{resolution}{Solução}
  \onslide<+->{No plano \(y=0\) a equação
  \[
    y^2 - x^2 = z
  \]}
  \onslide<+->{se reduz a}
  \begin{align*}
    \onslide<+->{0-x^2 & = z    \\}
    \onslide<+->{z     & = -x^2}
  \end{align*}
  \onslide<+->{%
    que corresponde uma parábola com vértice na origem\\
    e apontando para baixo
  }
\end{resolution}

%------------------------------------------------------------------------------%
\begin{resolution}{Solução}
\begin{columns}
\column{0.5\linewidth}
  \centering
  \(z=0\) \\
  \medskip
  \includegraphics{C_04_exemplo_3_z0}
\column{0.5\linewidth}
  \centering
  \(y=0\) \\
  \medskip
  \includegraphics{C_04_exemplo_3_y0}
\end{columns}
\end{resolution}
%------------------------------------------------------------------------------%

%------------------------------------------------------------------------------%
\begin{question}
  Considerando a esfera de raio $3$ centrada na origem
  \[
    x^2 + y^2 + z^2 = 9
  \]
  e o plano horizontal
  \[
    z = 2
  \]
  Determine a curva de interseção das duas superfícies no espaço
\end{question}

%------------------------------------------------------------------------------%
\begin{resolution}{Solução}
  A equação
  \[
    x^2 + y^2 + z^2 = 9
  \]
  \pause
  restrita ao plano $z=2$,
  \pause
  se reduz a
  \[
    \pause
    x^2 + y^2 + 2^2 = 9
  \]
  \[
    \pause
    x^2 + y^2 = 5
  \]
  \pause
  que corresponde a uma circunferência de raio $\sqrt{5}$ \\
  \pause
  contida no plano $z = 2$ \\
  \pause
  e centrada no ponto \((0, 0, 2)\)
\end{resolution}

%------------------------------------------------------------------------------%



%------------------------------------------------------------------------------%
\begin{question}
  Considerando a esfera
  \[
    x^2 + y^2 + z^2 = 9
  \]
  e o plano
  \[
    x + y + z = 0
  \]
  Determine a projeção da interseção das superfícies no plano $xy$
\end{question}

%------------------------------------------------------------------------------%
\begin{resolution}{Solução}
\begin{columns}[t]
\column{0.5\linewidth}
  \onslide<+->{Isolamos $z$ na equação do plano}
  \begin{gather*}
    \onslide<+->{x + y + z = 0 \\}
    \onslide<+->{z = -x - y}
  \end{gather*}
\column{0.5\linewidth}
  \onslide<+->{Substituímos na equação da esfera}
  \begin{align*}
    \onslide<+->{x^2 + y^2 + z^2             & = 9           \\}
    \onslide<+->{x^2 + y^2 + (-x - y)^2      & = 9           \\}
    \onslide<+->{x^2 + y^2 + x^2 + 2xy + y^2 & = 9           \\}
    \onslide<+->{2x^2 + 2y^2 + 2xy           & = 9           \\}
    \onslide<+->{x^2 + y^2 + xy              & = \frac{9}{2}}
  \end{align*}
  \onslide<+->{que é uma cônica rotacionada}
\end{columns}
\end{resolution}

%------------------------------------------------------------------------------%



%------------------------------------------------------------------------------%
\begin{question}
  Considere as quádricas
  \[
    x^2 + y^2 + z^2 = 9
    \hphantom{z = x^2 + y^2}
    \quad \text{(esfera)}
  \]
  \[
    z = x^2 + y^2
    \hphantom{x^2 + y^2 + z^2 = 9}
    \quad \text{(paraboloide circular)}
  \]
  Encontre a equação da curva de interseção

  \pause
  \bigskip
  Dica: use \(u = x^2 + y^2\)
\end{question}

%------------------------------------------------------------------------------%
\begin{resolution}{Solução}
  \onslide<+->{Substituímos a equação do paraboloide
  \[
    z = x^2 + y^2
  \]}
  \onslide<+->{na da esfera}
  \begin{align*}
    \onslide<+->{x^2 + y^2 + z^2                      & = 9 \\}
    \onslide<+->{x^2 + y^2 + \left(x^2 + y^2\right)^2 & = 9 \\}
    \onslide<+->{u + u^2                              & = 9 \\}
    \onslide<+->{u^2 + u - 9                          & = 0}
  \end{align*}
\end{resolution}

%------------------------------------------------------------------------------%
\begin{resolution}{Solução}
  \onslide<+->{Resolvendo \(u^2 + u - 9 = 0\)}
  \[
    \onslide<+->{\Delta \;=\; b^2 - 4ac}
    \onslide<+->{       \;=\; 1^2 -4 \times 1 \times (-9)}
    \onslide<+->{       \;=\; 37}
  \]

  \[
    \onslide<+->{u \;=\; \frac{-b\pm\sqrt{\Delta}}{2a}}
    \onslide<+->{  \;=\; \frac{-1\pm\sqrt{37}}{2}}
  \]

  \onslide<+->{Como $u\geq 0$}
  \[
    \onslide<+->{u = \frac{-1+\sqrt{37}}{2}}
  \]
\end{resolution}

%------------------------------------------------------------------------------%
\begin{resolution}{Solução}
  Então
  \[
    x^2 + y^2 = \frac{-1 + \sqrt{37}}{2}
  \]
  \pause
  Uma circunferência de raio \(\sqrt{\frac{-1 + \sqrt{37}}{2}}\)
  centrada em \((0, 0, 0)\)
\end{resolution}

%------------------------------------------------------------------------------%


%------------------------------------------------------------------------------%
\section{Lista Mínima}

%------------------------------------------------------------------------------%
\begin{frame}{Lista Mínima}
  Cálculo Vol. 2 do Thomas 12\textsuperscript{a} ed. -- Seção 12.6
  \begin{enumerate}
    \item Estudar o texto da seção
    \item Encontre a equação e esboce os gráficos dos cortes,
          nos planos $x=0$, $y=0$ e $z=0$, das quádricas dos exercícios: 1-12
  \end{enumerate}

  \vfill
  Atenção: A prova é baseada no livro, não nas apresentações
\end{frame}

%------------------------------------------------------------------------------%
\end{document}
%------------------------------------------------------------------------------%
